\documentclass[letterpaper,11pt]{article}

\usepackage[shortlabels]{enumitem}
\usepackage[margin=1in]{geometry}
\usepackage[most]{tcolorbox}

\begin{document}
\title{{\bf Unit 3: Basic Neuroanatomy Project} }
\author{Name: Shreya Kasireddy}

\date{4/24/25}
\maketitle

\section*{Short Answer}
\begin{enumerate}[a)]
\item Describe several advantages and disadvantages of biological computation with the brain compared to machine learning

\begin{tcolorbox}
The brain is much more resilient to errors, while many machine learning and other computer algorithms can be highly sensitive to even the smallest of errors. This is one important advantage of biological computation with the brain. Another advantage is that the brain is incredibly energy efficient. Computers and machine learning consume large amounts of energy to performs its computations. One more advantage of biological computation with the brain is that it facilitates creativity and intutition, which artificial intelligence systems are not able to mimic. However, machine learning algorithms allow impressive efficiency and speed in computation while naturally, the human brain may take longer. 
\end{tcolorbox}

\item Speculate what aspects of the architecture of the brain may cause these advantages or disadvantages, and similarly comment on aspects of machine learning’s architecture

\begin{tcolorbox}
Because of the brain's frontal and temporal lobe, the brain is able to display creativity and intuition through emotion and personality. Machine learning programs operate in a more binary manner and are thus not able to mimic the functions behind these emotions and more complex thinking. The occipital lobe of the brian, however, more closely follows the operations behind a machine learning program because it manages the processing and recognition of visual information which is similar to how a machine learning algorithm can "learn" and grow to recognize patterns and classify information. 
\end{tcolorbox}

\item Brainstorm some marvelous schemes for integrating advantages from both ways of computing. Draw, write, scribble etc… When you are done, do a quick google for your best ideas to see if anyone has researched or tried them already!

\begin{tcolorbox}
I think one important way computers can imitate the brain is by becoming more energy-efficient. This is a crucial point of inspiration from the brain because artificial intelligence is so energy-demanding. By optimizing the way certain algorithms are run, including reducing or cutting out unnecessary steps or pathways, we can improve the energy consumption of computer programs. While researching this, I learned that through optimized search and sorting algorithms, we can create more energy-efficient artificial intelligence models. This is becoming increasingly more important and widespread in the tech sector. 
\end{tcolorbox}
    
\end{enumerate}
\end{document}
